% 原书第 158—160 页；第 10.4 节、完整假设讨论与证明。
\section{C-TMLE 渐近线性的一般定理}
\label{sec:ch10-asymptotic-linearity}

下面的定理证明：如果存在某个选择器 $h_{1n}$，使 C-TMLE 具有渐近线性，那么我们提出的采用 $h_n$ 的 C-TMLE 也具有渐近线性。

\begin{bookresult}{定理}{10.1}
令 $\Psi:\mathcal M\to\mathbb R^d$ 在任意 $P\in\mathcal M$ 处路径可微，其有效影响曲线为
$D^*(P)=D^*(Q(P),G(P))$，并且 $\Psi(P)$ 只通过 $Q(P)$ 依赖于 $P$。为方便记号，我们也将 $\Psi(P)$ 记为 $\Psi(Q)$。假设
\[
 \Psi(Q)-\Psi(Q_0)
 =-P_0D^*(Q,G)+R_{20}(Q,G,Q_0,G_0),
\]
其中，余项 $R_{20}()$ 具有双重稳健结构，使得对所有 $Q\in\mathcal Q(\mathcal M)$，均有
$R_{20}(Q,G_0,Q_0,G_0)=0$。
令 $(G_{n,h}:h)$ 是 $G_0$ 的一族候选估计量，以标量 $h$ 为索引；令 $(Q_{n,h}^*:h)$ 是在靶向步骤中使用 $G_{n,h}$ 的一族 TMLE，从而对所有 $h$ 都有
$P_nD^*(Q_{n,h}^*,G_{n,h})=0$。
令 $h_n$ 是一个给定的选择器，并令 $Q_n^*=Q_{n,h_n}^*$ 为对应的 TMLE。令 $(Q^*,G_0)$ 为 $(Q_{n,h_n},G_{n,h_n})$ 的极限，使得当 $n\to\infty$ 时，
\[
 P_0\{D^*(Q_{n,h_n}^*,G_{n,h_n})-D^*(Q^*,G_0)\}^2
 \longrightarrow_p0.
\]
我们作如下假设。

\noindent\textbf{存在所需的选择器：}
假设对于一个固定的 $Q^*\in\{Q(P):P\in\mathcal M\}$，存在一个趋于零的序列 $h_{1n}$，满足
\begin{equation}
 \begin{aligned}
 &P_0D^*(Q^*,G_{n,h_{1n}})-P_0D^*(Q^*,G_0)\\
 &\qquad=(P_n-P_0)D_1(P_0)+o_P(n^{-1/2}).
 \end{aligned}
 \tag{10.11}
\end{equation}

\noindent\textbf{选择器求解关键方程：}
对于这样的选择器 $h_{1n}$，有
\begin{equation}
 \left\{\frac{d}{dh_n}P_nD^*(Q_n^*,G_{n,h_n})\right\}
 \ast(h_n-h_{1n})=o_P(n^{-1/2}),
 \qquad Q_n^*=Q_{n,h_n}^*.
 \tag{10.12}
\end{equation}

\noindent\textbf{二阶余项可以忽略：}
\[
 \begin{aligned}
 &P_0D^*(Q^*,G_{n,h_{1n}})-P_0D^*(Q^*,G_{n,h_n})\\
 &\quad-\left\{\frac{d}{dh_n}P_0D^*(Q^*,G_{n,h_n})\right\}
       \ast(h_{1n}-h_n)=o_P(n^{-1/2}),
 \end{aligned}
\]
\[
 (P_n-P_0)\{D^*(Q_{n,h_n}^*,G_{n,h_n})-D^*(Q^*,G_0)\}
 =o_P(n^{-1/2}),
\]
\[
 \left\{(P_n-P_0)\frac{d}{dh_n}D^*(Q^*,G_{n,h_n})\right\}
 \ast(h_{1n}-h_n)=o_P(n^{-1/2}),
\]
\[
 \begin{aligned}
 &P_n\left\{\frac{d}{dh_n}
   \{D^*(Q_n^*,G_{n,h_n})-D^*(Q^*,G_{n,h_n})\}\right\}\\
 &\qquad\ast(h_{1n}-h_n)=o_P(n^{-1/2}),
 \end{aligned}
\]
\[
 \begin{aligned}
 &\{P_0D^*(Q_n^*,G_{n,h_n})-P_0D^*(Q_n^*,G_0)\}\\
 &\quad-\{P_0D^*(Q^*,G_{n,h_n})-P_0D^*(Q^*,G_0)\}
 =o_P(n^{-1/2}).
 \end{aligned}
\]
那么，$\Psi(Q_n^*)$ 是 $\Psi(Q_0)$ 在 $P_0\in\mathcal M$ 处的渐近线性估计量，影响曲线为 $D^*(Q,G_0)+D_1(P_0)$：
% 原书前一句为 Q，下一式为 Q^*，按各自印刷保留。
\[
 \Psi(Q_n^*)-\Psi(Q_0)
 =(P_n-P_0)\{D^*(Q^*,G_0)+D_1(P_0)\}+o_P(n^{-1/2}).
\]
\end{bookresult}

\noindent\textbf{关于定理 10.1 假设的讨论。}
下面说明如何验证这些假设。第一个假设 (10.11) 要求存在序列 $h_{1n}$，使估计足够欠平滑，从而使光滑泛函 $\Phi_{Q^*}(G_n)$ 成为 $\Phi_{Q^*}(G_0)$ 的渐近线性估计量，其中
$\Phi_{Q^*}(G)=P_0D^*(Q^*,G)$。
当 $Q^*\ne Q_0$、即 $Q^*$ 设定错误时，这一点可能无法做到。此时，一个值得研究的问题是，能否推广下面的证明，说明即便如此，$\Psi(Q_n^*)$ 的二阶项表现仍然优于不求解关键方程 (10.12) 的标准 TMLE。
如果 $P_nD^+(Q_n^*,G_{n,h_n})=0$，第二个假设 (10.12) 就成立；但如前所述，我们不要求精确求解。我们的主张是，这个假设归结为要求
\[
 \left\{\frac{d}{dh_n}P_nD^*(Q_n^*,G_{n,h_n})\right\}
 \ast h_n=o_P(n^{-1/2}),
\]
因为可以将 $h_{1n}$ 选为仍能保持渐近线性 (10.11) 的最快趋零速率。
如果 $h\mapsto P_0D^*(Q^*,G_{n,h})$ 二次连续可微，第一个“二阶余项”假设就相当于要求
$(h_{1n}-h_n)^2=o_P(n^{-1/2})$。
只要 $h_{1n}^2$ 与 $h_n^2$ 都是 $o_P(n^{-1/2})$，这一条件就成立。因此，按照对 $h_n$ 的偏差解释，$G_{n,h_n}$ 的偏差必须以快于 $n^{-1/4}$ 的速率趋于零。
假设 $D^*(Q_n^*,G_{n,h_n})$ 以趋于 $1$ 的概率落在某个 $P_0$-Donsker 类中，$D^+(Q_n^*,G_{n,h_n})$ 也以趋于 $1$ 的概率落在某个 $P_0$-Donsker 类中，并且
\[
 P_0\{D^*(Q_n^*,G_{n,h_n})-D^*(Q^*,G_0)\}^2
 \longrightarrow_p0,
\]
\[
 P_0\{D^+(Q_n^*,G_{n,h_n})-D^*(Q^*,G_0)\}^2
 \longrightarrow_p0.
\]
% 原书后一式 D^+ 减去 D^*，并未印成 D^+ 减 D^+。
那么，由经验过程理论以及 $\max(h_{1n},h_n)\longrightarrow_p0$，可以推出第二、第三和第四个“二阶余项”假设成立。因此，这三个二阶余项假设只依赖于 $(Q_n^*,G_{n,h_n})$ 对 $(Q^*,G_0)$ 的一致性，以及一个 Donsker 类条件。最后，对于某些泛函 $H_1,H_2$ 和 $f$，第五个“二阶余项”假设中的余项一般可以表示为
\[
 \begin{aligned}
 \int &(H_1(Q_n^*)-H_1(Q^*))(H_2(G_{n,h_n})-H_2(G_0))\\
      &{}\times f(Q_n^*,G_n,Q^*,G_0)\,dP_0.
 \end{aligned}
\]
此时，如果排除奇点，即假设强正值性，那么只要
\[
 \|H_1(Q_n^*)-H_1(Q^*)\|_{P_0}\|H_2\|_{P_0}
 =o_P(n^{-1/2}),
\]
这个假设就成立。
% 原书此范数乘积中 H_2 未带两个 G 的差值，照印刷保留。

\begin{proof}[定理 10.1 的证明]
有
\[
 0=(P_n-P_0)D^*(Q_{n,h_n}^*,G_{n,h_n})
   +P_0D^*(Q_{n,h_n}^*,G_{n,h_n}).
\]
根据二阶假设，
\[
 (P_n-P_0)D^*(Q_{n,h_n}^*,G_{n,h_n})
 =(P_n-P_0)D^*(Q^*,G_0)+o_P(n^{-1/2}).
\]
令 $Q_n^*=Q_{n,h_n}^*$。由
\[
 \Psi(Q_n^*)-\Psi(Q_0)
 =-P_0D^*(Q_n^*,G_0)+R_{20}(Q_n^*,G_0,Q_0,G_0),
\]
以及假设 $R_{20}(Q_n^*,G_0,Q_0,G_0)=0$，可得
\[
 \begin{aligned}
 P_0D^*(Q_n^*,G_{n,h_n})
 ={}&P_0D^*(Q_n^*,G_0)\\
    &{}+\{P_0D^*(Q_n^*,G_{n,h_n})-P_0D^*(Q_n^*,G_0)\}\\
 ={}&\Psi(Q_0)-\Psi(Q_n^*)\\
    &{}+P_0D^*(Q_n^*,G_{n,h_n})-P_0D^*(Q_n^*,G_0).
 \end{aligned}
\]
因此，我们已经证明
% 原书此式在 D^*(Q^*,G_0) 与 P_0D^*(Q_n^*,G_{n,h_n}) 之间漏印 +；保留相邻乘写。
\[
 \begin{aligned}
 \Psi(Q_n^*)-\Psi(Q_0)
 ={}&(P_n-P_0)D^*(Q^*,G_0)P_0D^*(Q_n^*,G_{n,h_n})\\
    &{}-P_0D^*(Q_n^*,G_0)+o_P(n^{-1/2}).
 \end{aligned}
\]
由二阶假设，有
\[
 \begin{aligned}
 &P_0D^*(Q_n^*,G_{n,h_n})-P_0D^*(Q_n^*,G_0)\\
 &\qquad=P_0D^*(Q^*,G_{n,h_n})-P_0D^*(Q^*,G_0)
        +o_P(n^{-1/2}).
 \end{aligned}
\]
现在写成
\[
 \begin{aligned}
 &P_0D^*(Q^*,G_{n,h_n})-P_0D^*(Q^*,G_0)\\
 &\qquad=P_0D^*(Q^*,G_{n,h_n})-P_0D^*(Q^*,G_{n,h_{1n}})\\
 &\qquad\quad{}+P_0D^*(Q^*,G_{n,h_{1n}})-P_0D^*(Q^*,G_0).
 \end{aligned}
\]
根据“存在所需的选择器”假设，第二项等于
$(P_n-P_0)D_1(P_0)+o_P(n^{-1/2})$。
根据二阶假设，还可在 $h_n$ 处作如下 Taylor 展开：
\[
 \begin{aligned}
 &P_0D^*(Q^*,G_{n,h_n})-P_0D^*(Q^*,G_{n,h_{1n}})\\
 &\qquad=-\left\{\frac{d}{dh_n}P_0D^*(Q^*,G_{n,h_n})\right\}
          \ast(h_{1n}-h_n)+o_P(n^{-1/2}).
 \end{aligned}
\]
接着，写成
\[
 \begin{aligned}
 &-\left\{\frac{d}{dh_n}P_0D^*(Q^*,G_{n,h_n})\right\}
        \ast(h_{1n}-h_n)\\
 &\qquad=\left\{\frac{d}{dh_n}(P_n-P_0)D^*(Q^*,G_{n,h_n})\right\}
        \ast(h_{1n}-h_n)\\
 &\qquad\quad{}-P_n\left\{\frac{d}{dh_n}D^*(Q^*,G_{n,h_n})\right\}
        \ast(h_{1n}-h_n).
 \end{aligned}
\]
根据二阶假设，右端第一项为 $o_P(n^{-1/2})$。对于第二项，写成
\[
 \begin{aligned}
 &P_n\left\{\frac{d}{dh_n}D^*(Q^*,G_{n,h_n})\right\}
        \ast(h_{1n}-h_n)\\
 &\qquad=P_n\left\{\frac{d}{dh_n}D^*(Q_n^*,G_{n,h_n})\right\}
        \ast(h_{1n}-h_n)\\
 &\qquad\quad{}+P_n\left\{\frac{d}{dh_n}D^*(Q^*,G_{n,h_n})\right\}
        \ast(h_{1n}-h_n)\\
 &\qquad\quad{}-P_n\left\{\frac{d}{dh_n}D^*(Q_n^*,G_{n,h_n})\right\}
        \ast(h_{1n}-h_n).
 \end{aligned}
\]
根据二阶假设，右端第二项为 $o_P(n^{-1/2})$。由选择器 $h_n$ 的关键方程假设，右端第一项也是 $o_P(n^{-1/2})$。因此，我们已经证明
\[
 \Psi(Q_n^*)-\Psi(Q_0)
 =(P_n-P_0)\{D^*(Q^*,G_0)+D_1(P_0)\}+o_P(n^{-1/2}),
\]
定理得证。
\end{proof}

